\documentclass[a4paper]{article}     
\usepackage{amsfonts}
\usepackage{amsmath}
\usepackage{amssymb}
\usepackage{graphicx}

\newcommand{\asa}{\Leftrightarrow}
\newcommand{\adR}{\Rightarrow}
\newcommand{\adL}{\Leftarrow}
\newcommand{\Lh}{\left(}
\newcommand{\Rh}{\right)}
\newcommand{\Lv}{\left[}
\newcommand{\Rv}{\right]}
\newcommand{\La}{\left\{}
\newcommand{\Ra}{\right\}}
\newcommand{\Ras}{\right.}


\begin{document}

\noindent \large Auteur: Yoren Collier \\
\noindent \large Studierichting: Informatica\\
\noindent \large Oefeningenreeks 1E nr. 3\\

\medskip

\normalsize

$Z(1) \adR$ rest $5$\\

$Z(-1) \adR$ rest $1$\\

Doordat je de veelterm $Z$ deelt door een veelterm van de 1ste graad en een rest v/d 0-de graad krijgt, kun je afleiden dat de veelterm $Z$ van de 1ste graad is.\\

$\adR$ Als je de veelterm $Z$ deelt door een veelterm v/d 2de graad is de rest een veelterm v/d 1ste graad.\\

$Z(z) = D(z)\cdot Q(z) + R(z)$, met $Z$ de veelterm, $D$ de deler, $Q$ het quotiënt en $R$ de rest allemaal in functie van $z$.\\

$Z(z) = D(z)\cdot Q(z) + R(z) \asa Z(z) - D(z)\cdot Q(z) = R(z)$\\

$Z(1) = R(1)$ en $Z(-1) = R(-1)$\\

$R(z) = az + b$\\

$R(1) = a + b = 5$\\

$R(-1) = -a + b = 1$\\

$\left \{
\begin{array}{l} % dit is de l van links en niet het cijfer 1 !
a + b = 5 \\
-a + b = 1
\end{array}
\right .
\asa \left \{
\begin{array}{l}
a + 1 + a = 5\\
b = 1 + a
\end{array}
\right .
\asa \left \{
\begin{array}{l}
2a = 4 \asa a = 2\\
b = 3
\end{array}
\right .$\\

$\adR R(z) = 2z +3$\\

$Z(z) = (z + 1)(z - 1)\cdot Q(z) + 2z +3$


\begin{thebibliography}{999}
%\bibitem{a} Referentie
\end{thebibliography}
\end{document} 